\begin{description}
	\item[1行目] 一方のデータをオブジェクト\texttt{X}に代入しています。\texttt{c}は結合combineという関数です。
	\item[2行目] 他方のデータをオブジェクト\texttt{Y}に代入しています。
	\item[3行目] t検定を行っています。引数として，データ\texttt{X,Y}を指定し，次に\texttt{paired}オプションを\texttt{T}に\footnote{大文字の\texttt{T}は$TRUE$，真であり，この逆は\texttt{F}とかいて$FALSE$,偽を意味する\textbf{特別な文字}です。真か偽か，本当か嘘かという2つの状態を表す論理値と呼ばれるもので，ここではスイッチオン/オフの意味です。対応スイッチ，オン！という意味ですね。大文字の$T,F$をオブジェクト名にして利用することも可能ですが，Rが専門的な意味で使う予約語なので，一般的にはそのような使用は避けた方が望ましいです。}，対立仮説alternativeを「XがYよりも小さい」という\underline{片側検定}にしてあります。
\end{description}
